\section{Alternative priors and operational stability conditions}
\label{sec:alternative-priors}

This section gives two changes with complete statistical bridges. A soft
inverse-margin penalty retains the support of the raw Gaussian. A bounded
activation-fault condition weakens the previous hard support rule and has a
direct interpretation in terms of numerical errors and neuron failures.
Neither is a theorem for the \emph{unchanged} Gaussian prior.

\subsection{A soft prior with the same support as the raw Gaussian}

Let $\mu_{N,t}$ be the raw Gaussian of Section~\ref{sec:model}, $0<t\le1$,
and put $a=N/t^2\ge1$. Retain the postactivation disturbance convention
\[
 \|\eta\|_*^2=\sum_{s=1}^{L-1}\frac{\|\eta_s\|_2^2}{N}
                         +|\eta_{\rm out}|^2.
\]
Define the exact all-input distance to a zero output by
\[
 \rho(w)=\min_{x\in\{-1,1\}^d}
       \inf\{\|\eta\|_*:z_w(x;\eta)=0\},\qquad
 r(w)=\min\{1,\rho(w)\}.
\]
An output-only disturbance makes $\rho$ finite. If all clean outputs are
nonzero, continuity and finiteness of the input cube make $\rho>0$.
Every disturbance ball with radius strictly below $\rho$ preserves the clean
signs: a sign change would yield a zero along the segment from zero noise.
The function $\rho$ is measurable, since the noise balls are compact and
the input set is finite. Under the raw Gaussian, all clean outputs are
nonzero almost surely. To see this, condition successively on the preceding
layers: each nonzero state gives nondegenerate Gaussian preactivations,
and finally a nondegenerate Gaussian scalar readout. A finite union over
inputs completes the argument.

Fix constants $p>0$ and $\lambda>0$, and define
\begin{equation}\label{eq:soft-margin-prior}
 \Phi(w)=r(w)^{-p},\qquad
 d\nu_{N,t}(w)=Z^{-1}\exp[-a\lambda\Phi(w)]\,d\mu_{N,t}(w).
\end{equation}
On the null set $\rho=0$, set the density to zero. The normalizer is positive
and finite, and $\nu_{N,t}$ and $\mu_{N,t}$ are mutually absolutely
continuous. Thus no set of positive Gaussian probability is excluded.
In energy notation, the new density is proportional to
$\exp\{-a[\kappa_N(w)+2\lambda\Phi(w)]/2\}$.
One can optionally retain the original clean-margin restriction
$\min_x|z_w(x)|>1/2$; without that optional restriction, even the original
hard output-margin condition is unnecessary for the following result.
In both versions, the likelihood is noiseless \emph{sign consistency} with
the observed training labels.

\begin{theorem}[Soft robustness prior]
\label{thm:soft-margin-prior}
Suppose that, for a Boolean target $f$, there is a set $A_f$ of exact sign
representations such that
\[
 \mu_{N,t}(A_f)\ge e^{-aD},\qquad
 \rho(w)\ge r_*\quad(w\in A_f),
\]
where $D\ge1$ and $0<r_*\le1$ are independent of $N,t$.
For fixed $p,\lambda$, the consistency posterior under
\eqref{eq:soft-margin-prior} has a Gibbs-risk sample bound polynomial in
\[
 d,L,D,r_*^{-1},\varepsilon^{-1},\log(1/\delta).
\]
If the cylinder hypotheses hold at every $N\ge C$ and $0<t\le1$, the
guarantee holds on one sample event simultaneously for all those widths
and scales. The optional clean-margin version has the same conclusion
when $A_f$ also satisfies that clean margin.
\end{theorem}

\begin{proof}
Set $B=D+\lambda r_*^{-p}$. For every sample $S$ from $f$, its unnormalized
partition function satisfies
\[
 Z_S=\int \mathbf1_{\text{consistent with }S}
                 e^{-a\lambda\Phi}\,d\mu_{N,t}\ge e^{-aB}.
\]
The global prior normalizer cancels after conditioning. The posterior
$Q_S$ therefore obeys
\[
 D_{\rm KL}(Q_S\Vert\mu_{N,t})
  =-a\lambda\mathbb E_{Q_S}\Phi-\log Z_S\le aB,
 \qquad
 Q_S\{\Phi>u\}\le e^{a(B-\lambda u)}.
\]
Choose $0<\beta\le1/2$ and
$u=[B+\log(1/\beta)]/\lambda\ge1$. Since $a\ge1$, the displayed tail is
at most $\beta$. Let $Q_S'$ be $Q_S$ conditioned on $\Phi\le u$. Its
explicit density gives
\[
 D_{\rm KL}(Q_S'\Vert\mu_{N,t})
  \le aB+\log\frac1{1-\beta}\le a(B+1).
\]
Its support is strictly sign-robust at radius
$r_0=\tfrac12u^{-1/p}$. Apply the uniform robust posterior-compression
theorem in Appendix~\ref{sec:robust-entropy} with $K=B+1$ and $r=r_0$. It gives one
finite class $\mathcal H$, with polynomial log cardinality, and a common
candidate marginal whose prediction can be coupled to $Q_S'$ at each
input with disagreement at most $\alpha$.
Since $Q_S'$ is training-consistent, the finite-class relative Chernoff
bound from that report gives
\[
 R(Q_S)\le\beta+3\alpha+
       \frac{8[\log|\mathcal H|+\log(1/\delta)]}{n}
\]
with probability at least $1-\delta$. Choose $\beta,\alpha$ proportional
to $\varepsilon$. All class parameters are independent of $N,t$, which
also proves simultaneity. The same calculation with the extra clean-margin
indicator proves the optional version.
\end{proof}

The robust digital cylinder of Appendix~\ref{sec:cylinder} has
$D=\operatorname{poly}(d,L,C)$ and $r_*^{-1}=\operatorname{poly}(C)$.
Consequently Theorem~\ref{thm:soft-margin-prior} gives a polynomial-depth
bound for every $N\ge C$ and every $\gamma\ge1$ while retaining the support
of the raw Gaussian. The prior probabilities have changed. The result
therefore does not settle whether ordinary independent Gaussian weights
alone give the same guarantee.
An inverse-squared margin ($p=2$) is a particularly simple choice.
A penalty merely logarithmic in $1/\rho$ would only produce an exponentially
small retained radius through this proof, so would not give the same
polynomial conclusion.

\subsection{Robustness only within the activation range}

Call a disturbance \emph{bounded-state} if every perturbed hidden activation
belongs to $[-1,1]^N$. This allows replacing an activation by another legal
activation value, and leaves scalar output noise unrestricted. Require
sign stability only against bounded-state disturbances in the specified
RMS ball.

\begin{proposition}[Bounded-state compression]
\label{prop:bounded-state-compression}
The deterministic and posterior compression constructions in Appendix~\ref{sec:robust-entropy} remain valid under bounded-state robustness in place of unrestricted
postactivation robustness. The previously established target Gaussian-mass
estimates remain valid.
\end{proposition}

\begin{proof}
On every successful compression event, every virtual hidden state is a
$\tanh$ output. First-layer clipping occurs before the activation;
active-coordinate clipping, Gaussian replacement, and source reconstruction
also feed the activation. The exact postactivation errors comparing the
virtual computation to the original weights therefore produce states in
$[-1,1]^N$, and have the same total RMS bound as before. The margin argument
adds only output noise. Subsequent finite-grid rounding uses that positive
output margin and does not invoke robustness again. The small-width case
also uses only output noise. These are all uses of robustness in the two
proofs. Finally, the previous target cylinder satisfies the stronger,
unrestricted condition, so its probability lower bound is unchanged.
\end{proof}

This relaxation can be substantial at a fixed radius. With one hidden
layer, let $n=N-1$, $c=0.9$, and take a constant hidden state
$h=(c,\ldots,c,-c)$ and readout
$v=(2/(nc),\ldots,2/(nc),1/c)$. Its output is $1$.
Unrestricted hidden and output noise can erase it at norm
$[N\|v\|^2+1]^{-1/2}=O(N^{-1/2})$.
For bounded-state noise, the last coordinate can decrease by at most
$1-c$. The remaining coordinates and scalar output noise can decrease the
output by at most
$\sqrt{4N/(nc^2)+1}\,\|\eta\|_*$.
Thus a zero requires
\[
 \|\eta\|_*\ge\frac{2c-1}{\sqrt{4N/n+c^2}}>0.26
 \qquad(N\ge2).
\]
This separates the notions of robustness; its readout energy grows with
$N$, so it is not a sample lower bound or a separation of the calibrated
prior supports.

\subsection{Small numerical errors and rare neuron faults}

For $a>0$, $0\le p\le1$, and $b\ge0$, say a network tolerates $(a,p,b)$ if it preserves its own clean signs at
every input under all simultaneous bounded-state disturbances such that,
in each hidden layer, at most $pN$ coordinates have $|\eta_i|>a$, and
$|\eta_{\rm out}|\le b$. Exceptional coordinates may be replaced
arbitrarily within $[-1,1]$. This combines numerical errors at all neurons
with an adversarial fraction of neuron faults. Dropout and pruning are
particular faults; average performance under random dropout is a weaker
property than this adversarial condition.

\begin{lemma}[Fault and RMS bridges]
\label{lem:fault-rms}
Tolerance of $(a,p,b)$ implies bounded-state RMS robustness at
$r=\min\{a\sqrt p,b\}$. Conversely, unrestricted RMS robustness at $\tau$
implies tolerance of $(a,p,b)$ whenever
\[
 (L-1)(a+2\sqrt p)^2+b^2\le\tau^2.
\]
\end{lemma}
\begin{proof}
If the total RMS error is at most $r$, each layer has at most
$Nr^2/a^2\le pN$ coordinates with error larger than $a$; output noise is
at most $r\le b$. Conversely, the ordinary coordinates contribute at most
$a\sqrt N$ to a layer's Euclidean error. Each faulty coordinate differs
by at most $2$, so their contribution is at most $2\sqrt{pN}$.
Sum the squared layer RMS bounds and the scalar output bound.
\end{proof}

For a single rule independent of target and data, retain the original
clean margin and define
\begin{align*}
 E_v&=N\|v\|^2,&
 \tau(w)&=\frac{0.01}{\sqrt{1+(L-1)E_v}},\\
 a(w)&=\frac{\tau(w)}{4\sqrt L},&
 p(w)&=\frac{\tau(w)^2}{64L},&
 b(w)&=\frac{\tau(w)}2.
\end{align*}
Restrict the Gaussian to networks tolerating these faults, then condition
on noiseless training sign consistency. Every network in the previous
radius-$\tau(w)$ support belongs to this support: the second bound in
Lemma~\ref{lem:fault-rms} is at most $\tau(w)^2/2$.
In particular, the previous target cylinder still lies in the support.

\begin{corollary}[Polynomial-depth learning with activation faults]
\label{cor:fault-prior}
For this fixed fault-tolerance support rule, the same qualitative
polynomial-depth sample theorem holds for all $\gamma\ge1$ and every
width above the layered Boolean target circuit's width.
\end{corollary}
\begin{proof}
The same target partition bound and Gaussian readout tail as before allow
conditioning on $E_v\le B$ while discarding posterior mass at most
$\beta$, where $B=\operatorname{poly}(d,L,C,\log(1/\beta))$.
Every retained network has bounded-state robustness radius at least
\[
 r_B=\frac{10^{-4}}{32L[1+(L-1)B]},
\]
because $a(w)\sqrt{p(w)}=\tau(w)^2/(32L)\le b(w)$.
The conditional posterior has the same polynomial scaled-KL bound as in
the previous argument. Apply Proposition~\ref{prop:bounded-state-compression}
and the finite-class generalization step. All parameters are independent
of width and Gaussian scale, giving the simultaneous conclusion.
\end{proof}

This condition has an operational interpretation, but it is not established
as a universal outcome of ordinary gradient descent. Likewise, random-noise
stability alone cannot simply replace it. A network that reads a single
hidden coordinate can be stable under small independent per-coordinate
Gaussian noise while that coordinate can be erased with only
$O(N^{-1/2})$ RMS adversarial noise. Testing particular dropout masks or
one quantization grid similarly does not control the disturbances used by
the compression coupling.
The conservative constants may give $p(w)N<1$ at small widths; then the
allowed number of arbitrary faults is zero, and the condition concerns
bounded errors at all coordinates. No claim of tolerance to a practically
large dropout fraction follows from these constants.
